\documentclass[12pt, a4paper]{article} % Dokumentenklasse

\usepackage[utf8]{inputenc} % Zeichensatz und Schrift
\usepackage[T1]{fontenc}
\usepackage{lmodern} % nice alternative is {mathpazo}
\usepackage[english]{babel}

\usepackage[left=2.5cm, right=2.5cm, top=2.5cm, bottom=2.50cm]{geometry} % Seitenformat

\usepackage{amsmath} % Mathematikzeichen
\usepackage{amsfonts}
\usepackage{amssymb}
\usepackage{mathtools}

\usepackage{booktabs} % Tabellen
\usepackage{array}
\usepackage{dcolumn}
\usepackage{tabularx}
\usepackage{threeparttable}

\usepackage{graphicx} % Grafiken
\usepackage{subfigure}



\usepackage{lipsum} % Paket für Random-Text (nur für Anschauungszwecke benötigt, kann gelöscht werden)

\usepackage{setspace} % Zeilenabstand
\onehalfspacing

\usepackage{placeins}

%% Literaturverzeichnis und Zitation % %

\usepackage[%
citestyle=authoryear-comp,%
bibstyle=JME,%
maxbibnames=99,% %maximum number of names printed in bibliography before truncation with ``et al.'' is used
maxcitenames=2,
datezeros=false,% no leading 0 if dates are printed
date=long,%
isbn=false,% show no ISBNs
natbib=true,% enable natbib-compatibility
url=false,% show no urls
backend=biber % use Biber here
]{biblatex}

% vergrößert Abstand zwischen den Quellen im Lit.verzeichnis

\setlength\bibitemsep{1.5\itemsep}

% setzt "et al." bei deutscher Sprache

\DefineBibliographyStrings{ngerman}{%
  andothers = {et\addabbrvspace al\adddot}
}


% vermeidet Witwen und Waisen

\clubpenalty10000 
\widowpenalty10000
\displaywidowpenalty=10000 



% Hier .bib-Datei einfügen

\bibliography{bspbib}

\usepackage[font=scriptsize, labelfont=bf]{caption}

%---------------------------------------------------------------------------------------
%
%								HIER BEGINNT DAS DOKUMENT
%
%---------------------------------------------------------------------------------------


\begin{document}

% TITELSEITE

\begin{titlepage}
\begin{center}
\includegraphics[scale=0.65]{KITlogo}\\
%\Large
%Karlsruher Institut für Technologie\\
\large
Fakultät für Wirtschaftswissenschaften\\
Institut für Volkswirtschaftslehre (ECON)
\vfill{
%\\\vspace{2cm}
\Large
Bachelor Thesis in Macroeconomics\\\vspace{0.5cm}
,,US commercial banks``\\\vspace{0.5cm}
Winter semester 2019/20% \vspace{2cm}
}
\vfill{
\LARGE
% -------------------------------------------------------------------------------------
% Ihr Thema
Not yet defined
% -------------------------------------------------------------------------------------
\\ 
\vspace{0.5cm}
\normalsize
% -------------------------------------------------------------------------------------
% Nummer Ihres Themas
(Topic 1)
% -------------------------------------------------------------------------------------
}
\end{center}

\vfill{
\normalsize
% -------------------------------------------------------------------------------------
% Ihre pers\"{o}nliche Daten
\noindent Alexander Schlechter \\
Matr.-Nr. 2054108 \\
alexander.schlechter@student.kit.edu}
% -------------------------------------------------------------------------------------
\end{titlepage}
\newpage

% VERZEICHNISSE

\pagenumbering{Roman}
\tableofcontents
% \listoffigures
% \listoftables
\newpage
\pagenumbering{arabic}


% ================================================================================
% 									Hauptteil
% ================================================================================


\section{Introduction}

Explorative journey through financial filings of us commercial banks.

Different perspectives: 
	- banking sector as whole (Aggregated)
	- Median banks

Analysis timeseries and cross section

\section{Main part}

\subsection{Data}

Data from  $http://pages.stern.nyu.edu/~pschnabl/data/data_callreport.htm$ ,originally from FIEC031 Reports.
Schneider adjust for variables change and so on

Timeframe: 1976Q1-2013Q4, data per year and quarter for each individual bank

We take our crisis data from the official documentation of the National Bureau of Economic Research(NBER. One can also differentiate between a so called "banking" (originated in the banking sector) vs "market" (originated from outside banking sector) crisis as in Berger and Bouwman (2013).  The assumption here is that banking crisis could be reflected stronger in the data. The crisis will be marked on the graphs as red zones.
Crisis:


\begin{itemize}
\item 1980Q1-1980Q3 early 1980s recession (market crisis)
\item 1981Q3-1982Q4 early 1980s recession (market crisis)
\item 1990Q3-1991Q2 credit crunch (banking crisis)
\item 2001Q2-2001Q4 dotcom bubble (market crisis)
\item 2007Q4-2009Q3 financial crisis (banking crisis)
\end{itemize}

Following the convention of the Federal Reserve Bulletin, banks are allocated by size into following categories:
\begin{itemize}
 \item 10 largest banks
 \item large banks (those ranked
  11 through 100)
 \item medium-sized banks (those ranked 101 through
 1,000)
 \item small banks (those ranked 1,001 and higher)
 \end{itemize} 

Other important structural events to consider:
\begin{itemize}
\item Repeal of Glass-Steagall act 
\end{itemize}

negative equity removed

Investment banks becoming commercial banks in 2009 removed

bank mergers removed by growth of assets above $40\%$

for risk analysis: commercial banks aggregated to bhcs

Values are in thousands

\subsection{Methods}


No inflation adjustment

Tranformations: Log

Timeseries analysis:

Detrending with HP Filter 1600

Deseasonalize with X11 procedure

Correlation/Autocorrelation

Significance



\subsection{General look at us banks}

The following two sections divided into assets and liabilities will guide us through the distribution of financial components held by the us commercial banks sector as a whole. We will see what the types and amounts of financial instruments banks are holding and how they evolved over time. 

\subsubsection{Assets Liabilities}

The following section will show some insights about the asset side of the whole us commercial banking sector. We want to describe and analyse how it evolved over time and find anomalies 
The typical positions a commercial bank is holding on the asset side of its balance sheet are:

\begin{itemize}
 \item Cash
 \item Fed funds sold and securities purchased under agreements to resell
 \item Securities
  	\begin{itemize}
 	\item Treasury
 	\item Mortage backed security
 	\item Other
 	\end{itemize}
 \item Loans net\footnote{Loans and leases net of unearned income and allowance for loan and lease losses}
 \item Trading assets
 	\begin{itemize}
 	\item Interest rate derivatives
 	\item Other fixed income
 	\item Other trading
 	\end{itemize}
 \item Other assets \footnote{composed of derivatives "not for trading" and other items}
 \end{itemize} 
 
 Every position beside the trading assets are hold "not for trading purposes". Meaning for instance the securities position and loans position are held until maturity. 
 Figure \ref{fig:assets} gives a general overview how the total assets held by all banks per year and quarter evolved over time. The value of assets rose from below 2 trillion to above 13 trillions dollars of assets. In comparison, the GDP of US in 2013 was at 16.78 trillions. We can also see a period of flatness/low growth from  year 1990-1993, a drop in 2002 and 2008/Q2. While the absolute values give us some indications of where anomalies occur, we will perform time-series analysis by decomposing the logarithmised data in trend, cyclical components. As seen within Figure \ref{fig:assets}, the value of assets has a clear trend with no indication of existing seasonality over the given timeframe. The second graph within Figure \ref{fig:assets} compares the development of assets transformed by natural logarithm to the absolute values. Both are drawn within their own vertical axis. While the growth of absolute assets is more exponential over time, the logarithm of assets grows more linear. This linear trend tells us that the relative growth rate stays constant over time. In the third part of Figure \ref{fig:assets} we can see the de-trended part of the asset timeseries. Here we can see variations that significantly differ from the trend. As mentioned in the data section the red areas show crisis periods. The alignment of crisis periods with us banks total assets variations is limited. We can see that the impact of the early 1980s recession did not lead to more volatility than other normal periods. The dotcom bubble in 2001 lead to a downward variation of us commercial banks assets away from the trend. In regards to the financial crisis in 2008 we see a huge positive variation with a rapid drop back to the trend. Assuming we have set the most fitting filtering parameters for the HP-Filter, it is interesting to see that the assets did not fall significantly below the trend. The loss was rather caused by an overheated market falling back to normal.
 
  
\begin{figure}[hbtp]
\centering
\caption{Asset side}
\includegraphics[scale=0.3]{gr\\

Diving into more detailed analysis of the balance sheet positions, Figure \ref{fig:positions} gives us the detrended development of each individual position for both the asset(left column) and liabilities side(right column).
The cash position shows some clear spike in 2008, beside then the movements show no clear sign of irregularity. With context in mind it makes sense that banks liquidated some of their assets to increase cash. The rise in cash comes along with a significant fall in securities. Figure \ref{fig:sharepositions} showing the share of the positions further supports our thesis of sold securities to raise cash. While the share of securities fell until 2008, share of cash rose.
Trading assets follow the same behaviour as total assets in crisis 2008, but its variation in crisis periods do not significantly differ to other periods.  
Other interesting observations on the asset side, is a significant spike in

The volatility of fedfundsrepoasset increased from 1996 onwards. 
fedfundsrepoliab drops in 2002
The zero values for trading assets and liabilities are either caused because they are not reported or were not existent. 

\begin{figure}[hbtp]
\centering
\caption{Detrended asset positions(left column) \hspace{2cm} Detrended liabilities positions(right column)}
\includegraphics[scale=0.3]{graphs/DescriptiveStats/OtherAnalysis_clean_PositionsCyclical_7613.png}
\label{fig:positions}
\end{figure}


\begin{figure}[hbtp]
\centering
\caption{Share of asset positions - unstacked}
\includegraphics[scale=0.3]{graphs/DescriptiveStats/OtherAnalysis_ShareofAssetsPlot_7613.png}
\label{fig:sharepositions}
\end{figure}


\noindent \textit{Graph description}: Figure 1 shows the aggregates of the main variables from the asset side of the balance sheet over time. Figure 2 shows the share of each aggregated balance sheet position of all commercial banks over time. Figure 3 plots the share of each balance sheet position unstacked.\\


\noindent \textit{Key Observations:}
\begin{itemize}
\item loans make up the largest share of assets
\item share of trading assets have risen as well as interest rate derivatives
\item loans and trading assets have risen more than securities in timeframe year 2000-2009
\item Share of trading assets peaked in 2008 while securities fell.
\item There is a noticeable anomaly in year 2002. Significant amounts of repo lending is transferred into other assets. Other assets are derivatives not available for sale. 
\item drop in assets in 2002 and 2008
\item share in cash has fallen until 2008 and then increased again
\item Share of cash continously fell until 2008, and then it increased significantly again
\end{itemize}

\newpage

\begin{figure}[hbtp]
\centering
\caption{Correlation of balance sheet positions (assets side)}
\includegraphics[scale=0.5]{graphs/DescriptiveStats/OtherAnalysis_AssetsCorrelationHeatMap_7613.png}
\end{figure}

\noindent \textit{Graph description}: Figure 3 shows the Pearson Correlation Coefficient between different balance sheet positions\\

\textit{Key Observations:}
\begin{itemize}
\item Cash has the weakest correlation with the rest positions
\item No significant stand outs
\end{itemize}

\newpage

\subsubsection{Liabilities}

\if(false)
\begin{figure}[hbtp]
\centering
\caption{Liabilities side}
\includegraphics[scale=0.3]{graphs/DescriptiveStats/OtherAnalysis_LiabilitiesDistribution_7613}
\end{figure}




\begin{figure}[hbtp]
\centering
\caption{Share of liabilities positions}
\includegraphics[scale=0.3]{graphs/DescriptiveStats/OtherAnalysis_ShareofLiabilities_7613}
\end{figure}

\newpage

\begin{figure}[hbtp]
\centering
\caption{Share of liabilities positions}
\includegraphics[scale=0.3]{graphs/DescriptiveStats/OtherAnalysis_ShareofLiabilitiesPlot_7613}
\end{figure}
\fi

\begin{figure}[hbtp]
\centering
\caption{Share of liabilities positions}
\includegraphics[scale=0.3]{graphs/DescriptiveStats/OtherAnalysis_ShareofLiabilitiesAll3_7613.png}
\end{figure}


\noindent \textit{Graph description}: The graph shows the aggregates of the main variables from the liabilities side of the balance sheet over time.\\

\noindent \textit{Key Observations:}
\begin{itemize}
\item Deposits as main source of funding
\item Irregularities in year 2002: repos drop while other liabilities rise
\item Irregularity in year 1983 might be caused by measuring/reporting differences
\item In 2008 share of deposits at lowest point. Although the aggregated assets peaked at that time. 
\item Deposits continuously decreased from 1990 onwards. Other financing such as "Other borrowed money" and "trading liabilities" rose
\item repos share decreased significantly until end of 2013
\end{itemize}



\begin{figure}[hbtp]
\centering
\caption{Correlation of balance sheet positions (liabilities side)}
\includegraphics[scale=0.5]{graphs/DescriptiveStats/OtherAnalysis_LiabilitiesCorrelationHeatMap_7613}
\end{figure}

\newpage

\subsection{Growth}


\begin{figure}[hbtp]
\centering
\caption{Growth of assets}
\includegraphics[scale=0.3]{graphs/DescriptiveStats/OtherAnalysis_GrowthAssetsYear_7613}
\end{figure}

\noindent \textit{Graph description}: The graph shows the annual growth rate of aggregated assets of all commercial banks. Two investment banks, who did become commercial banks in 2009, are excluded. \\

\noindent \textit{Key Observations:}
\begin{itemize}
\item Three negative growth rates in year 1991, 2001, 2010
\end{itemize}

\begin{figure}[hbtp]
\centering
\caption{Growth of top 1 percent banks assets}
\includegraphics[scale=0.3]{graphs/DescriptiveStats/OtherAnalysis_GrowthAssetsYear1_7613.png}
\end{figure}

\noindent \textit{Graph description}: The graph shows the annual growth rate of aggregated assets of top top $1\%$ commercial banks. Two investment banks, who did become commercial banks in 2009, are excluded. \\

\noindent \textit{Key Observations:}
\begin{itemize}
\item More negative growth rates in 1990 and 2010
\item No negative growth in 2001

\end{itemize}

\begin{figure}[hbtp]
\centering
\caption{Growth of all banks vs top 1 percent}
\includegraphics[scale=0.3]{graphs/DescriptiveStats/OtherAnalysis_GrowthTrends_7613}
\end{figure}

\noindent \textit{Graph description}: Annual growth rate with trend for all banks and top $1\%$. The second graph shows the cycle part from the time series filter. \\

\noindent \textit{Key Observations:}
\begin{itemize}
\item Top $1\%$ growth rates are more volatile
\item Pearson Correlation between all banks vs top $1\%$, Significance: (0.6371607133788253, 1.696703469447756e-05)
\item 1986, 2001, 2006 almost 0,02 difference toward trend
\end{itemize}

\newpage


\begin{figure}[hbtp]
\centering
\caption{Banks default}
\includegraphics[scale=0.3]{graphs/DescriptiveStats/NewLeverage_BanksDefaultPercentage_7613.png}
\end{figure}

\noindent \textit{Graph description}: The graph shows an estimation of how many banks have defaulted at a certain time (year,quarter). It is based on the negative equity recorded by banks. Hence, it is not exact and some banks might continue to exist in case of mergers or bailouts. Also sometimes banks are double counted, if a negative equity does not immediately result in bankruptcy. \\
\noindent \textit{Key Observations:}
\begin{itemize}
\item main defaults in years 1986-1991 and 2009-2011
\item long stable period from 1991-2008
\item In 1990 there were many more smaller banks. Smaller banks might have a higher likelihood to fail. In 1990:  $~74\%$ small banks, 2010: $~35\%$ small banks 
\end{itemize}


\subsection{Loans}

\begin{figure}[hbtp]
\centering
\caption{Loans}
\includegraphics[scale=0.3]{graphs/DescriptiveStats/OtherAnalysis_LoanDistribution_7613}
\end{figure}

\noindent \textit{Graph description}: It shows the share of loan types of total loans over time.\\

\noindent \textit{Key Observations:}
\begin{itemize}
\item real estate loans has largest share

\end{itemize}

\begin{figure}[hbtp]
\centering
\caption{Loans by repricing maturity}
\includegraphics[scale=0.3]{graphs/DescriptiveStats/OtherAnalysis_ShareofMaturityLoans_9713}
\end{figure}


\begin{figure}[hbtp]
\centering
\caption{Residential Loans by repricing maturity}
\includegraphics[scale=0.3]{graphs/DescriptiveStats/OtherAnalysis_ShareofMaturityResLoans_9713}
\end{figure}

\newpage

\subsection{Distribution of asset sizes among banks - To Big to Fail}

\begin{figure}[hbtp]
\centering
\caption{Count of banks by percentiles}
\includegraphics[scale=0.8]{graphs/DescriptiveStats/BankCounts}
\end{figure}

\newpage

\begin{figure}[hbtp]
\centering
\caption{Aggregate assets by percentiles}
\includegraphics[scale=0.3]{graphs/DescriptiveStats/OtherAnalysis_ShareofAssetsByBanksPercentiles_7613.png}
\end{figure}

\iffalse
\begin{figure}[hbtp]
\centering
\caption{Top 10 banks assets vs rest}
\includegraphics[scale=0.3]{graphs/DescriptiveStats/LeverageDistribution_AssetsTop10vsRest_7813}
\end{figure}

\begin{figure}[hbtp]
\centering
\caption{Rise of top 10 banks asset share}
\includegraphics[scale=0.3]{graphs/DescriptiveStats/LeverageDistribution_shareTop10_7813}
\end{figure}
\fi


\subsection{Median banks by asset size}

In Figure 15, we have in the left column the asset interval size and in the corresponding row the number of banks per year.\\

\begin{figure}[hbtp]
\centering
\caption{Banks count by asset size}
\includegraphics[scale=1]{graphs/MedianBank_FrequencyTable_8010}
\end{figure}


\textbf{Typical small/medium/large bank}

Banks are assigned three different buckets (small/medium/large) depending on asset size.\\

Small bank: $0<$ assets $<= 10^5$

Medium bank: $10^5<$ assets $<= 10^6$

Large bank: $10^6<$ assets 

\begin{figure}[hbtp]
\centering
\caption{Asset size by bank}
\includegraphics[scale=0.3]{graphs/DescriptiveStats/MedianBank_SmallvsMiddlevsBigMedianBankAssets_7613.png}
\end{figure}

\begin{figure}[hbtp]
\centering
\caption{Medium vs large bank by asset size}
\includegraphics[scale=0.3]{graphs/DescriptiveStats/MedianBank_MiddlevsBigMedianBankAssets_7613.png}
\end{figure}


\begin{figure}[hbtp]
\centering
\caption{Small bank: liability side}
\includegraphics[scale=0.3]{graphs/DescriptiveStats/MedianBank_ShareofLiabilitiesSmallBank_7613.png}
\end{figure}


\begin{figure}[hbtp]
\centering
\caption{Medium bank: liability side}
\includegraphics[scale=0.3]{graphs/DescriptiveStats/MedianBank_ShareofLiabilitiesMediumBank_7613.png}
\end{figure}


\begin{figure}[hbtp]
\centering
\caption{Large bank: liability side}
\includegraphics[scale=0.3]{graphs/DescriptiveStats/MedianBank_ShareofLiabilitiesBigBank_7613.png}
\end{figure}



\newpage

\subsection{Leverage}

Throughout the analysis the definition of accounting leverage (assets/equity) is used. Equity is calculated by total assets minus total liabilities. In addition, for risk analysis banks belonging a bank holding company were aggregated. Hence, the dataset which was used contained bhcs and independent banks. Figure 21 below shows the mean, median and weighted-average leverage by asset size for each point in time. The weighted-average leverage ratio is calculated by taking into account the asset size for each bank every point in time. Every leverage ratio for each individual bank is only accounted in the weighted-average by its share of assets compared to the total assets of all banks at that point of time.

\begin{figure}[hbtp]
\centering
\caption{Median and Average leverage for all banks}
\includegraphics[scale=0.3]{graphs/NewLeverage_LeverageRatioAllBanks_7613}
\end{figure}

\newpage

\begin{figure}[hbtp]
\centering
\caption{Average/Mean leverage plots}
\includegraphics[scale=0.3]{graphs/NewLeverage_LeverageRatioMean_7613}
\end{figure}


\begin{figure}[hbtp]
\centering
\caption{Median leverage plots}
\includegraphics[scale=0.3]{graphs/NewLeverage_LeverageRatioMedian_7613}
\end{figure}

\newpage

\begin{figure}[hbtp]
\centering
\caption{Trend plots}
\includegraphics[scale=0.3]{graphs/NewLeverage_LeverageRatioTrends_7613.png}
\end{figure}


\begin{figure}[hbtp]
\centering
\caption{Cyclical plots}
\includegraphics[scale=0.3]{graphs/NewLeverage_LeverageRatioCyclical_7613.png}
\end{figure}

\newpage

\begin{figure}[hbtp]
\centering
\caption{Correlation cyclical leverage}
\includegraphics[scale=1]{graphs/DescriptiveStats/Correlation_Cyclical_Leverage.png}
\end{figure}

\noindent \textit{Graph description}: Figures 21-25 focus on leverage ratios for every year and quarter over all banks. Banks with equity or assets below zero are excluded.


\noindent \textit{Key Observations:}
\begin{itemize}
\item Extreme outliers of leverage in year 1992/93 and 2009 lead to spikes in average leverage.
\item Average-weighted leverage significantly higher from 1976 to 2004. Aligns with higher leverage of top $0.1\%$ and $1\%$ banks. Hence, the overall risk from year 1976-2004, was driven high by the larger banks. 
\item Figure 21: All banks median leverage, seasonal effect every year?
\item Leverage lowest in 2007
\item Overall Leverage did fall over time: Introduction of Basel 1 in 1988 might have lead to continuously decrease in leverage
\item Top $0.1\%$ have much higher volatility, which could just be caused by the low sample size.
\item Top $0.1\%$ and $1\%$ actually become less risky than all banks together from 2010 onwards
\item Figure 25 shows slightly negative correlation between the top $0.1\%$ banks and all banks. 
\item Crisis 2008/9: Top $0.1\%$ banks have a average leverage spike in 2008, while top $10\%$, $50\%$ and all  spike around a year later.  
\end{itemize}

\newpage

\begin{figure}[hbtp]
\centering
\caption{Boxplots (1976-2013)}
\includegraphics[scale=0.3]{graphs/Leverage/LeverageDistribution_LeverageRatioBoxplot_7613.png}
\end{figure}

\noindent \textit{Graph description}: Boxplots of all leverage ratios by banks by year. 
\\

\noindent \textit{Key Observations}:
\begin{itemize}
\item $75\%$ of all banks have a leverage ratio between 10-15.
\end{itemize}



\iffalse
\begin{figure}[hbtp]
\centering
\caption{Leverage Top 10 vs Rest over all years}
\includegraphics[scale=0.4]{graphs/Leverage/LeverageDistribution_LeverageRatioTop10vsRest_7813}
\end{figure}


\newpage

\noindent \textit{Graph description}: Since the top 10 banks share of assets did rise up to 60\% in 2013, it is important to differentiate. The graph shows the average leverage (assets/equity) for every year quarter 4.\\ Blue Line: Top 10 Banks by assets\\
Orange Line: All banks beside the top 10\\

\noindent \textit{Key Observations}:
\begin{itemize}
\item Leverage of top 10 banks tends to be higher
\item Trend of falling leverage is similar
\end{itemize}




\begin{figure}[hbtp]
\centering
\caption{Leverage Top10 vs Rest detailed look into crisis}
\includegraphics[scale=0.4]{graphs/Leverage/LeverageDistribution1_LeverageDetailCrisisTop10vsRest_0510}
\end{figure}


\noindent \textit{Key Observations}:
\begin{itemize}
\item Top10 banks leverage peak in year 2008/3 before the rest banks leverage peak in year 2009/4 (theory of risky assets from big banks to small transfer?) 
\end{itemize}
\fi
\pagebreak

\textbf{A look into the distribution of leverage}\\



\begin{figure}[hbtp]
\centering
\caption{Distribution 1980-2013}
\includegraphics[scale=0.3]{graphs/Leverage/LeverageDistribution_LeverageRatio_8013}
\end{figure}

\noindent \textit{Graph description}: Counts are normed to 1. Leverage is transformed with log10. Leverage ratios are always from quarter 4.
\\

\noindent \textit{Key Observations}:
\begin{itemize}
\item Log-normal distribution
\item Large standard deviation in year 2010 with 18.82
\item Less and higher bars in 2012 indicate higher homogeneity in 2013 compared to the years before.
\end{itemize}

\pagebreak


\begin{figure}[hbtp]
\centering
\caption{Distribution in crisis 2003-2011}
\includegraphics[scale=0.3]{graphs/Leverage/LeverageDistribution_LeverageRatio_0311}
\end{figure}

\noindent \textit{Graph description}: Leverage ratios are always from quarter 4 and  logarithmised.

\noindent \textit{Key Observations}:
\begin{itemize}
\item Increasing homogenity over time.
\end{itemize}

\section{Evaluation/Outlook}

shadow banking/investment banks not considered

off balance sheet

valuations not realistic, book values...

Applying more time-series models

Applying of models such as Regression...

Other structural events considered

significant part of trading assets still owned by non commercial banks

\section{Conclusion}




% ================================================================================
% 							Literaturverzeichnis
% ================================================================================

\newpage
\printbibliography[
heading=bibintoc,
title={Bibliography}
]

crisis data: https://www.nber.org/cycles.html


% ================================================================================
% 								Appendix
% ================================================================================

%\appendix
%\section{Anhang}




\end{document}